\section{Introduction}
\label{sec:introduction}

Randomized controlled trials (RCTs) are the cornerstone of scientific discovery across medicine, social science, and public policy~\cite{ImbensRu15}.
While the statistical foundations of experimentation are well-developed, the question of how experimental designs affect downstream decision-making and outcomes is not as well-understood.
This is especially true when treatments are evaluated in a series of experiments, as is common in both natural and social sciences~\cite{ChenWi13,KasySa21}.

While there is a rich literature on single-stage designs, experimentation is rarely a one-off process: across domains ranging from drug development (FDA trial stages) to educational interventions (IES goals or project types), a common structure is that a smaller RCT determines whether a treatment merits further investigation, followed by a larger-scale confirmatory trial that informs adoption decisions~\cite{FDACD22,DepartmentofEducationBo22}.
This pilot to scale-up to adoption paradigm is society's solution to a fundamental resource allocation problem---how to efficiently screen promising interventions while managing costs and risks.
The role of small RCTs in screening promising interventions is critical: both rejecting an effective intervention and scaling up a harmful one can be costly to society~\cite{RossellMu07}.
Correspondingly, funding and regulatory bodies have highlighted the need for methodological progress on the design of early-stage trials~\cite{Tipton21,FDAHH20,DepartmentofEducationBo22}.

In practice, due to the tight resource constraints of most pilot studies, researchers often rely on \emph{convenience samples}---participants recruited from subpopulations where studies are easiest to conduct~\cite{TiptonSp21}.
When the effectiveness of the treatment under study varies with subject characteristics, the treatment effect in such a non-representative sample may be significantly better or worse than in the target population.
A common proposal to mitigate these errors is to recruit \emph{representative samples} that better match the target population~\cite{FDAHH20,DepartmentofEducationBo22}.
This is in line with a broad push for generalizability or external validity in science~\cite{EgamiHa23, BryanTi21, Rothwell05, HenrichHe10, VoelklVo18}.

In this work, we study how the \emph{design of a pilot\footnote{We focus on the use of small RCTs as screening mechanisms for scale-up decisions, and refer to these as \emph{pilot studies}.
Pilots have many definitions across fields, and sometimes refer to even earlier stages of experimentation where the goal is ``not to test the effectiveness of the intervention but to assess the feasibility of the protocol'' \cite[Figure 2]{EldridgeLa16}.} RCT} affects the outcomes of a target population when the pilot determines whether an intervention merits further study.
We characterize how the impact of pilot design propagates through three stages: the pilot study, the follow-up study, and the adoption decision, as shown in \Cref{fig:exp_process_flow}.
By modeling the multi-stage nature of scientific discovery, we are able to analyze the role statistical significance tests play as primary continuation criteria for making decisions in the current scientific status quo~\cite{FDA19,Berry91}.
In doing so, we are neither arguing for significance tests as an optimal continuation criterion, nor that impact is a normatively correct goal of any given stakeholder.
Instead, we seek to understand how a researcher whose work is evaluated by a significance test might design studies if they care about impact, and how pushing for more representative samples without changing how studies are evaluated might affect research outcomes.

\begin{figure}[ht]
  \centering
  \vspace{0.8\baselineskip}
  \begin{tikzpicture}[scale=0.19]
  \tikzstyle{every node}+=[inner sep=0pt]
  \pgfmathsetmacro{\diagArrowLabelGap}{1.8}
  \newcommand{\dashedDecisionArrowLabel}[3]{%
    \path let
    \p1 = (#1),
    \p2 = (#2),
    \n1 = {atan2(\y2-\y1,\x2-\x1)}
    in node[
      rotate=\n1,
      text=red,
      fill=white,
      inner sep=1pt
    ] at ($($(#1)!0.5!(#2)$)+(\n1+90:\diagArrowLabelGap)$) {#3};%
  }

  \draw (14.2,-20.3) node {\textbf{Pilot RCT}};
  \draw (33.1,-20.3) node {Follow-up RCT};
  \draw (52.3,-20.3) node {Adoption};
  \draw (14.2,-31.9) node[text=red] {Statistical Test};
  \draw (33.1,-31.9) node[text=red] {Statistical Test};
  \draw (52.3,-31.9) node {Treatment Effect};

  \draw [black, dashed] (16.87,-29.54) -- (30.43,-22.66);
  \dashedDecisionArrowLabel{16.87,-29.54}{30.43,-22.66}{Continue?}
  \fill [black] (30.43,-22.66) -- (29.49,-22.58) -- (29.94,-23.47);
  \draw [black, dashed] (35.78,-29.56) -- (49.62,-22.64);
  \dashedDecisionArrowLabel{35.78,-29.56}{49.62,-22.64}{Adopt?}
  \fill [black] (49.62,-22.64) -- (48.68,-22.55) -- (49.12,-23.45);

  \draw [black] (14.2,-22.66) -- (14.2,-29.54);
  \fill [black] (14.2,-29.54) -- (14.7,-28.74) -- (13.7,-28.74);

  \draw [black] (33.1,-22.64) -- (33.1,-29.56);
  \fill [black] (33.1,-29.56) -- (33.6,-28.76) -- (32.6,-28.76);

  \draw [black] (52.3,-22.64) -- (52.3,-29.56);
  \fill [black] (52.3,-29.56) -- (52.8,-28.76) -- (51.8,-28.76);

\end{tikzpicture}

  \caption{
    \textbf{The three-stage experiment pipeline we model.}
    The intervention is first evaluated in a budget-constrained pilot RCT, where the result of a significance test determines whether it advances to a follow-up RCT.
    A second significance test determines adoption of the intervention in the target population, leading to a potential improvement in average outcomes if successful.
  }
  \label{fig:exp_process_flow}
\end{figure}

We show that when optimizing for expected downstream impact, neither \emph{representative sampling} nor \emph{convenience sampling} is optimal across all regimes.
In fact, \textbf{the optimal pilot design depends heavily on the budget of the pilot study}, as illustrated in \Cref{fig:small-budget-vs-proportional-vs-feasible}.
Under mild assumptions, we characterize the optimal pilot design for sufficiently small and asymptotically large pilot budgets under \emph{general prior distributions} of the conditional average treatment effects across subpopulations.
In the large-budget regime where pilot studies achieve high statistical precision, our model verifies the conventional wisdom that pilot designs that match the composition of the target population achieve optimal expected downstream outcomes.
However, when the budget constraint is severe, the optimal pilot design takes the opposite extreme to a representative sample: it concentrates sampling on a single subpopulation.
This provides surprising evidence that the convenience samples commonly used in early-stage experiments may have close to optimal impact, and that representative samples may be counterproductive for highly constrained pilot studies when statistical significance is important.
Because our objective yields dramatically different designs under small budgets than those found in work on external validity~\cite{Tipton13,EgamiLe23,Bouyamourn25} and minimax-regret objectives~\cite{HuZh24,OleaPr25}, we believe substantial work remains to understand resource-constrained experiment design for decision-making.

\begin{figure}[ht]
  \vspace{-0.2\baselineskip}
  \centering
  \includegraphics[width=0.73\textwidth]{figures/pilot_impact_small_budget_vs_large_budget_vs_feasible.pdf}
  \vspace{-0.2\baselineskip}
  \caption{
    \textbf{Optimal design depends on pilot resources.}
    We plot the expected downstream impact achieved by the optimal pilot design, the best small-budget single-subpopulation design, and a representative sample.
    In budget-constrained settings, optimal pilots sample from a \emph{single homogeneous subpopulation}, while representative sampling is best for well-resourced trials.
  }
  \label{fig:small-budget-vs-proportional-vs-feasible}
  \vspace{-0.2\baselineskip}
\end{figure}

To gain insight into how treatment effect heterogeneity---the key motivation for representative samples---drives optimal pilot design, we specialize our small-budget optimal design to elliptical priors, which generalize multivariate normal distributions.
Here, we show that the single subpopulation that is optimal to sample under any sufficiently small budget is given by an index trading off
\begin{enumerate}[label=(\roman*)]
  \item a subpopulation's expected effect size,
  \item how predictive its effect is of the effect in the follow-up and target populations, and
  \item its recruitment cost.
\end{enumerate}

Why are homogeneous samples optimal under small budgets, while under large budgets the optimal pilot design approaches representative sampling?
\emph{For small budgets, the optimal design is driven by the signal-to-noise ratio of the test, while for large budgets it is driven by the mismatch between the pilot estimand (the sample average treatment effect) and the downstream estimand (the population average treatment effect).}

Our theoretical analysis of small-budget optimal design applies under mild assumptions to a much broader class of settings than pilot studies.
Homogeneous (single-subpopulation) sampling is small-budget optimal when significance testing is used to decide whether to receive \emph{any} downstream payoff, so long as that payoff does not depend directly on the pilot outcomes.
We show this general result via an asymptotic expansion of the objective that shows that for sufficiently small budgets, it is approximately equivalent to maximizing a signal that is linear in the subpopulation sample sizes divided by a noise scale equal to the square root of the total sample size.
This signal-to-noise ratio behaves roughly like a quasiconvex function (its positive part is quasiconvex), and therefore under a linear budget constraint it is maximized at an extreme point of the feasible design polytope.

In the other extreme where the pilot budget is very large, the pilot test can be made arbitrarily precise, and a representative sample will make almost no type I or type II errors.
Any non-representative sample, on the other hand, will make errors on some non-negligible fraction of the prior distribution of treatment effects---those for which the signs of the sample and population average treatment effects differ.

We conclude this section with a heuristic description of what kinds of subpopulations are optimal to sample under a tight budget.
Consider our main objective of maximizing the expected impact of a pilot used to determine whether to run a follow-up RCT, and compare it to two other potential objectives: maximizing the impact of adopting based on just a single RCT, or maximizing the probability of passing the pilot significance test.
Under a sufficiently tight budget, our analysis shows that sampling from a single subpopulation is optimal for all three objectives.
However, which population to sample from varies.

Consider the case when all subpopulations have the same sampling costs, and a potential follow-up RCT will use a very large and representative sample.
In this case, pilot false positives never make it to adoption---the follow-up RCT filters out any treatments with negative effects in the target population.
Since false positives are not penalized, the pilot can be more aggressive in sampling from subpopulations that have larger expected effects, even if they are less predictive of the target population effect.
However, it cannot focus only on passing the pilot, as there is still an important trade-off between passing the pilot and avoiding false negatives:  avoiding rejecting interventions with a positive effect in the target population.
Thus, pilot impact falls between single-RCT impact and pilot pass probability in terms of how it prioritizes how promising a subpopulation is versus how predictive it is of the target population effect.

The remainder of the paper is organized as follows.
\Cref{sec:Model} presents a model of the experimental pipeline, including the treatment effect structure, experimental design decisions, and significance testing protocols, and formalizes the optimization problem.
Sections \ref{sec:large-budget} and \ref{sec:small-budget} present our main theoretical results characterizing optimal pilot design in the large-budget and small-budget regimes.
\Cref{sec:case-study} presents a semi-synthetic case study calibrated to a real set of experiments in education research, demonstrating that real-world pilot studies can fall squarely within either of the budget regimes we characterize.
Finally, we situate our contributions in the literature on experimentation and external validity in \Cref{sec:related-work}, and discuss implications for future work in \Cref{sec:conclusion}.
